\section{Bosonization and the horizontal gas}\label{c7:boson}

The conditional angular trace from Section~\ref{c7:contour} will now be
rewritten as a gas on the horizontal circle.  This change of description
has two analytic stages.  We first evaluate each fixed current coefficient
and prove that the resulting series is absolutely integrable under a real
tilted Gaussian law on contours inside the field's analytic strip.  After this integration,
the resulting meromorphic terms can cross the collision divisor; its residue
creates one additional particle type.  The exact gas identity is then proved
at each fixed site count for the sufficiently large integer periods covered
by the stability estimate.  The later limit in the period uses the original
finite residue regulator, rather than a termwise limit of the gas.

The charged-fermion and oscillator descriptions below implement the
boson--fermion correspondence~\cite{Frenkel1981}; a primary treatment
connecting Clifford operators and Heisenberg fields is
\cite[Sections~6.4 and 8.1]{FrenkelJingWang2000}.  We include the mode calculation
because the charge shifts and phases determine the scalar normalization.
Throughout the conditional calculations, $T>0$ and the analytic loop
coefficients are fixed, and $\mathbf q=e^{-2\pi/T}$.

\subsection{Charge sectors and the seam trace}

Work first with one species.  Write
$\Psi_j^+=\Psi_j$ and $\Psi_j^-=\Psi_{j^*}$ for $j=1,2,3$.
The eigenvalues $c=(c_1,c_2,c_3)\in\mathbb Z^3$ of the three zero modes
$h_{j,0}$ label charge sectors.  In each sector the nonzero modes satisfy
\[
[h_{j,n},h_{l,s}]=n\delta_{n,-s}\delta_{jl},\qquad
h_{j,n}^*=h_{j,-n},
\]
and generate an oscillator space from a ground state $|c\rangle$ of energy
$|c|^2/2$.  To define this ground state, for each $j$ fill the creators of
sign $\operatorname{sgn}c_j$ up to energy $|c_j|-1/2$, with the largest
energy leftmost, and put the three lists in the order $j=1,2,3$.
The anticommutation relations give the displayed commutator: changing
normal order contributes exactly $n$ mode pairs for $n>0$.  The positive
oscillator modes kill $|c\rangle$.

These oscillator spaces account for all fermionic states.  For one value of
$j$, encode a configuration by occupied positions in $\tfrac12+\mathbb Z$,
starting with the negative positions filled.  A positive-sign creator inserts
a particle at a positive position, and a negative-sign creator deletes one at
a negative position.  In charge $c_j$, the occupied positions in decreasing
order are
\[
c_j+\tfrac12-k+\lambda_k,\qquad k\ge1,
\]
for a unique partition $(\lambda_k)$.  Every partition occurs, and its
excess energy is $\sum_k\lambda_k$.  These are exactly the multiplicities
generated by the negative oscillator modes.

Let $R_j$ raise $c_j$ by one, commute with nonzero modes, and act on the
charge basis with Klein sign $(-1)^{\sum_{l<j}c_l}$.  In matrix elements the
field is
\begin{equation}\label{c7:vertex-field}
\Psi_j^\epsilon(v)=y^{1/2}R_j^\epsilon y^{\epsilon h_{j,0}}
 \exp\!\left(\epsilon\sum_{n>0}y^n h_{j,-n}/n\right)
 \exp\!\left(-\epsilon\sum_{n>0}y^{-n}h_{j,n}/n\right),
\qquad y=e^{2\pi i v/T},\quad \epsilon=\pm1.
\end{equation}
Both sides have the same commutators with all nonzero $h$ modes.  Their
matrix elements between ground states agree by inserting or deleting the
largest filled mode, which fixes the remaining factor and sign.  Powers of
$y$ use the indicated coordinate before exponentiation.

Recall from the preceding section that
\[
b=b_w e^{\mathsf A},\qquad b_w=b_0y^{-\ell},\qquad
b_0=e^{-i(\theta+\varphi_0)},\qquad
\mathsf A=m a_T-i\chi,
\]
where $\mathsf A$ has zero spatial mean, and that
\[
W_b=W_{b_w}\exp\!\left(\sum_n\mathsf A_n(-h_{1,n}+h_{2,n})\right).
\]
We also retain $Z_\Gamma=\mathrm{tr}\Gamma$ from the conditional trace.
The action of $W_{b_w}$ on a ground state is
\begin{equation}\label{c7:weyl-ground}
W_{b_w}|c\rangle
=(-1)^{c_1+c_1c_2}b_0^{\ell+c_2-c_1}
  |c_2+\ell,c_1-\ell,c_3\rangle,
\end{equation}
and it swaps oscillator labels 1 and 2.  Here is a verification that also
fixes its scalar.  In the three-root word for $W_{b_w}$, the modes with
$r+s=\ell$ form commuting blocks $(1+F)(1+E')(1+F)$, where
\[
F=b_0^{-1}\Psi_{4,r}\Psi_{2,s},\qquad
E'=-b_0\Psi_{5,-s}\Psi_{1,-r}.
\]
When $r,s$ have opposite signs the block fixes the vacuum.  When both are
positive it contributes $E'$, and when both are negative it contributes
$F$.  Put the species-1 creator first in each contributing block.  The sign
for grouping creators by species cancels the sign for reversing the second
energy list, giving
$W_{b_w}|0\rangle=b_0^\ell|\ell,-\ell,0\rangle$.
This calculation can be made on finite occupation sets, since only finitely
many blocks act nontrivially there.  Conjugation of fields sends
$\Psi_1\mapsto-b_0^{-1}y^\ell\Psi_2$ and
$\Psi_2\mapsto b_0y^{-\ell}\Psi_1$.  With the Klein signs, increasing
$c_1$ changes the ground coefficient by $(-1)^{1+c_2}b_0^{-1}$ and
increasing $c_2$ changes it by $(-1)^{c_1}b_0$.  Increasing $c_3$ changes
nothing.  These ratios and the vacuum value prove
\eqref{c7:weyl-ground}; the evenness of $\ell\in8\mathbb Z$ is used here.

Rotation sends a ground charge $d$ to $(d_2,d_3,-d_1)$ with sign
$(-1)^{\sum d_j+d_1(d_2+d_3)}$.  The first part counts the minus sign on
each creator, and the second moves the first species list past the other
two.  Applying rotation to \eqref{c7:weyl-ground} therefore gives
\begin{equation}\label{c7:seam-ground}
\Gamma W_{b_w}|c\rangle
=\mathbf q^{|c''|^2/2}(-1)^{c_2+c_3+c_2c_3}
 b_0^{\ell+c_2-c_1}|c''\rangle,
\qquad c''=(c_1-\ell,c_3,-c_2-\ell).
\end{equation}
On creation oscillators of frequency $n>0$, in addition to this ground
state action, it acts by
\begin{equation}\label{c7:oscillator-twist}
B_n=\mathbf q^n\operatorname{diag}(1,R),\qquad
R=\begin{pmatrix}0&1\\-1&0\end{pmatrix}.
\end{equation}

Consider a fixed coefficient in the three current exponentials on the two
sides of $W_b$.  By the cyclic rule \eqref{c7:cyclic-current}, its outgoing
field pairs can be put to the right of $W_b$ in the order $U,L,D$.
Here $U$ denotes the rightmost root, $L$ the cycled leftmost root, and $D$
the middle exponential; their individual labels will be specified below.
Each $\Psi$ carries again the physical multiplier
$(-2\pi i/(a_*T))^{1/2}$.  Consequently a current $X_{bc}(f)$ contributes
$(a_*/(2\pi i))f(v)\,dv$ together with its pair of physical fields.
For a field of sign $\epsilon$ at folded coordinate $v$, define its centered
coordinate $w$ by
\begin{equation}\label{c7:centered-fields}
w=v-i/2\quad(j=1,\text{ short}),\qquad
w=v-i\quad(j=2,\text{ outer}),\qquad
w=v\quad(j=3,\text{ outer}).
\end{equation}

Let the total inserted charges of types $1,2,3$ be $(S_1,E_2,F_3)$, and let
$n=(n_1,n_2,n_3)$ be the ground charge to the right of the fields.  Applying
\eqref{c7:seam-ground} to $c=n+(S_1,E_2,F_3)$ and imposing trace closure
gives
\begin{equation}\label{c7:charge-closure}
S_1=\ell,\qquad
n_2=(F_3-E_2-\ell)/2,\qquad
n_3=(-F_3-E_2-\ell)/2.
\end{equation}
The last two values are integral: every current has two fields, so
$S_1+E_2+F_3$ is even, and $\ell$ is even.  The charge $n_1$ remains free.
The $b_0$ factor becomes $b_0^{(F_3+E_2-\ell)/2-n_1}$.

For later sign accounting, order all fields from left to right.  If $N_j$ is
the number of fields of type $j$, the Klein exponent apart from inversions
is $n_1(N_2+N_3)+n_2N_3$.  Add the exponent in
\eqref{c7:seam-ground} and substitute \eqref{c7:charge-closure}.  The result
modulo two is
\begin{equation}\label{c7:klein-sign}
(F_3+E_2)/2+\#(3\text{ before }2)
                 +\#(\text{outer before short}).
\end{equation}
For example, $N_2+N_3$ is even and the non-inversion remainder reduces to
$n_2+F_3\equiv(E_2+F_3)/2\pmod2$.  The initial ground-coordinate power on
field $i$ of type $j$ is
$y_i^{1/2+\epsilon_i(n_j+\sum_{l>i,\,j_l=j}\epsilon_l)}$.

\subsection{Oscillator contractions at fixed insertion order}

For normalized creators with coefficient vector $u$ followed by normalized
annihilators with coefficient vector $v'$, the oscillator trace with
propagator $B_n$, divided by its empty trace, is
\[
\exp\!\left({v'}^t B_n(I-B_n)^{-1}u\right).
\]
For one oscillator with multiplier $t$, cyclicity and the commutator give
$\langle(a^*)^l a^l\rangle=l!(t/(1-t))^l$ and zero for unbalanced moments;
diagonalizing the unitary part gives the vector formula.  Reordering two
exponentials adds the contraction of the earlier annihilator with the later
creator.  An unsplit exponential adds half of its internal contraction.

Apply these rules to the seam source
$\exp(\sum_n\mathsf A_n(-h_{1,n}+h_{2,n}))$.  For $k=2\pi n/T>0$ and
$s=(-1,1,0)^t$, \eqref{c7:oscillator-twist} gives
\[
s^t\left(\tfrac12+B_n(I-B_n)^{-1}\right)s
=\tfrac12\coth(k/2)+\tfrac12\tanh k
=\frac{\cosh(3k/2)}{2\sinh(k/2)\cosh k}=E(k).
\]
Since $D(k)=kE(k)/(2\pi)$, the source self-contraction is
$\exp(\tfrac12\int_T\mathsf A D\mathsf A)$.  Its cross contractions dress
a positive short and a positive outer field respectively by
\begin{equation}\label{c7:oscillator-dresses}
f_s^{\rm osc}(w)=-\sum_{k=2\pi n/T\ne0}
 \frac{\mathsf A_n e^{ikw}}{2\sinh(k/2)},\qquad
f_o^{\rm osc}(w)=\sum_{k=2\pi n/T\ne0}
 \frac{\mathsf A_n e^{ikw}}{2\cosh k}.
\end{equation}
A negative field receives the negative exponent.  To check the centering in
this formula, the creator contraction uses $s^t(I-B_n)^{-1}$ and the
annihilator contraction uses $B_n(I-B_n)^{-1}s$.  The short coordinate has
$y^n=\mathbf q^{n/2}e^{ikw}$; the two outer coordinates have
$y^n=\mathbf q^ne^{ikw}$ and $e^{ikw}$.  Substitution gives the two displayed
multipliers.  The empty oscillator trace is
\begin{equation}\label{c7:empty-oscillator}
D_W=\prod_{n\ge1}(1-\mathbf q^n)^{-1}(1+\mathbf q^{2n})^{-1}.
\end{equation}

For fields $i<l$ put $d'=w_l-w_i$ and $Z=e^{2\pi i d'/T}$.  Each short field
has self-factor $H_s=\prod_{n\ge1}(1-\mathbf q^n)$, and its pair base is
\[
F_s(Z)=(1-Z)\prod_{r\ge1}(1-\mathbf q^rZ)(1-\mathbf q^r/Z).
\]
Each outer field has self-factor
$H_o=\prod_{r\ge1}(1-\mathbf q^{2r})^{(-1)^r}$, and its pair base is
\[
F_o(Z)=(1-Z)\prod_{r\ge1}
 [(1-\mathbf q^{2r}Z)(1-\mathbf q^{2r}/Z)]^{(-1)^r}.
\]
Every pair base is raised to $\epsilon_i\epsilon_l$.  Use $F_o(Z^{-1})$
when a type-3 field precedes a type-2 field.  For example, the outer block of
$B_n(I-B_n)^{-1}$ at $t=\mathbf q^n$ is
$(tR-t^2I)/(1+t^2)$; combining it with the two centered $y$ powers gives
this orientation.  There is no oscillator pair contraction between a short
and an outer field.

These formulas hold for the fixed current coefficient under discussion.
Here are the convergence details needed for that assertion.  Start with
outgoing folded coordinates in strict strip order
$0<\eta_1<\cdots<\eta_a<1$, where $\eta_i=\Im v_i$.  The uncycled pairs
start slightly above zero; the cycled pairs start slightly below one.
Cyclicity may first be applied to individual mode pairs.  At fixed analytic
loop coefficients, a quadratic Fourier mode of energy increment $\Delta N$
has only $O(1+N+|\Delta N|)$ potentially acting terms on energy $N$.
Exponential Fourier decay and the graded kernel estimate of
Lemma~\ref{c7:angular-trace} therefore make the integrated sums absolute
under $\mathbf q^{L_0}$.  Cycling reindexes these absolutely convergent
energy and occupation sums.

At separated field coordinates, let $N_0$ be the energy just left of the
first field, $N_i$ the energy to the right of field $i$, and $N_a$ the
terminal energy.  A field mode $r$ changes energy by $-r$, so its coordinate
power together with $\Gamma$ supplies the damping
\[
\exp\!\left[-\frac{2\pi}{T}\left(
 \eta_1N_0+\sum_{i=1}^{a-1}(\eta_{i+1}-\eta_i)N_i
 +(1-\eta_a)N_a\right)\right].
\]
The energy dimensions and the remaining graded kernels grow only
subexponentially.  Inside one bilinear no internal contraction is possible;
if its two coordinates coincide, the intermediate energy is bounded by the
sum of its bracketing energies and the number of mode splittings is
polynomial.  The same formula is therefore integrable on the original nearby
lines, also by a coincident-coordinate limit within a bilinear.

One may equivalently sum the same finite-level matrix elements in oscillator
bases and first truncate the oscillator modes.  After replacing coefficients
by their absolute values, $B_n$ is bounded by $\mathbf q^n$ times a
permutation.  Direct creator--annihilator sums have the ordered-radius gap,
thermal wraps have the remaining annulus gap, and source terms have the
Fourier decay of $\mathsf A$.  Their large-$n$ logarithms are summable.
The ground sums have Gaussian decay at this fixed field list and winding.
This permits removal of the truncations and proves the oscillator formulas
on the initial lines.  Subsequent fixed-term moves use analytic continuation
of these formulas.  Absolute integration over the Gaussian samples will be
proved separately below.

\subsection{Changing the period}

The product formulas that turn the angular contractions into horizontal
ones are
\begin{equation}\label{c7:modular-products}
\begin{split}
Z^{-1/2}F_s(Z)/H_s^2
&=\frac{-i}{T}e^{-\pi d'^2/T}
 \frac{2\sinh\pi d'\prod_{r\ge1}
 (1-e^{-2\pi(rT+d')})(1-e^{-2\pi(rT-d')})}{H_*^2},\\
H_s&=\sqrt T\,e^{-\pi T/12+\pi/(12T)}H_*,
\qquad H_*=\prod_{r\ge1}(1-e^{-2\pi Tr}),\\
Z^{-1/2}F_o(Z)
&=-i\mathbf q^{-1/2}\tanh(\pi d'/4)
 \prod_{r\ge1}\tanh(\pi(rT+d')/4)\tanh(\pi(rT-d')/4),\\
H_o&=\sqrt{T/8}\,\mathbf q^{-1/4}
 \prod_{r\ge1}\tanh(\pi rT/4).
\end{split}
\end{equation}
For the first identity, the two sides have the same simple zeros and
multipliers under $d'\mapsto d'+T,d'+i$.  Their quotient is a pole-free
doubly periodic function, and the derivatives at zero fix it.  The scalar
formula for $H_s$ can be obtained in the same comparison: after removing the
zero, compare the quadratic logarithmic coefficients.  With
$U(p)=\sum_{r\ge1}p^r/(1-p^r)^2=\sum_{r\ge1}rp^r/(1-p^r)$, divisor
interchange gives
\[
-\frac{\pi^2}{6T^2}+\frac{4\pi^2}{T^2}U(\mathbf q)
=\frac{\pi^2}{6}-\frac\pi T-4\pi^2U(e^{-2\pi T}).
\]
Multiplying this equality by $-1/(2\pi)$ gives
$\partial_T\log(H_s/H_*)=1/(2T)-\pi/12-\pi/(12T^2)$;
equality at $T=1$ fixes the scalar constant.  For the outer identity, compare
divisors and multipliers under $+T,+4i$.  The product of the left side and
its $2i$ translate is $-\mathbf q^{-1}$, fixing the magnitude of the
constant; a small positive real $d'$ fixes its sign.  The derivative at zero
then gives the formula for $H_o$.

Write $\alpha_p=\theta+\varphi_0+2\pi p$ and $g_w=2\pi\ell/T$.  For the
short fields let $N_s$ be their count and $X_s=\sum_i\epsilon_iw_i$.
Extracting $Z^{1/2}$ from their pair bases changes the coordinate powers to
$y_i^{\epsilon_i(n_1+\ell/2)}$.  The identities
\[
2\sum_{i<l}\epsilon_i\epsilon_l=\ell^2-N_s,
\qquad
\sum_{i<l}\epsilon_i\epsilon_l(w_l-w_i)^2
=\ell\sum_i\epsilon_iw_i^2-X_s^2
\]
show that the net $H_s$ power is $\ell^2$ and describe the pair Gaussian.
Completing the heat and coordinate powers leaves
$\mathbf q^{(n_1+\ell/2)^2/2+\ell^2/8}$ and
$e^{-i\alpha_p\ell/2}$.  Since $\ell/2$ is integral, Poisson summation gives
\[
\sum_{n_1\in\mathbb Z}
\mathbf q^{(n_1+\ell/2)^2/2}
 e^{i(\theta+\varphi_0)(n_1+\ell/2)}
 e^{2\pi i(n_1+\ell/2)X_s/T}
=\sqrt T\sum_{p\in\mathbb Z}
 e^{-T\alpha_p^2/(4\pi)-\alpha_pX_s-\pi X_s^2/T}.
\]
The $X_s^2$ term cancels the pair Gaussian.  Including physical field
normalizations, the short contribution apart from
$\sqrt T e^{-T\alpha_p^2/(4\pi)}$ is as follows:
\[
\begin{gathered}
\text{field dress: }\quad
\epsilon_i(-\alpha_pw_i-g_ww_i^2/2),\\
\text{scalar: }\quad
\exp[-\pi T\ell^2/12-\pi\ell^2/(6T)-i\alpha_p\ell/2],\\
\text{per-field factor: }\quad (2\pi/a_*)^{1/2}H_*.
\end{gathered}
\]
The remaining short pair base is $2\sinh\pi d'$ with its two image factors
from \eqref{c7:modular-products}.  The possible phase
$(-i)^{\ell^2/2}$ is one for $\ell\in8\mathbb Z$.

For the outer fields put $Y_i=e^{2\pi iw_i/T}$.  Extract the inverse square
root when type 3 precedes type 2 and the ordinary square root otherwise.
The residual power of every $Y_i$ is $-\ell\epsilon_i/2$.  The ground and
coordinate exponent of $\mathbf q$ before the modular factors is
\[
\frac{n_2^2+n_3^2}{2}+E_2n_2+\frac{E_2^2}{2}
=\frac{(E_2+F_3)^2+\ell^2}{4}.
\]
The modular pair and self factors remove $(E_2+F_3)^2/4$.  Orient every
remaining odd outer base as $d'=w_l-w_i$; this contributes the
3-before-2 sign.  The outer contribution is
\[
\begin{gathered}
\text{field dress: }\quad -i\epsilon_i(\alpha_p+g_ww_i)/2,\\
\text{scalar and phase: }\quad
\mathbf q^{\ell^2/4}e^{i\alpha_p\ell/2}
(-i)^{(E_2+F_3)^2/2},\\
\text{per-field factor: }\quad
(\pi/(4a_*))^{1/2}\prod_{r\ge1}\tanh(\pi rT/4).
\end{gathered}
\]
The remaining pair base is $\tanh(\pi d'/4)$ with its images.
Writing $E_2+F_3=2h$, the displayed phase is $(-1)^h$.  It cancels the
charge sign in \eqref{c7:klein-sign}, and the reorientation cancels its
3-before-2 sign.  The sign left is exactly
$(-1)^{\#(\text{outer before short})}$.  In combination with the odd
like-family bases, it makes the coordinate expression alternating under
exchange of any neighboring fields.  Whole even current pairs may therefore
be reordered without a sign.  The total winding scalar per species is
\begin{equation}\label{c7:modular-winding}
\exp[-\pi T\ell^2/12-2\pi\ell^2/(3T)].
\end{equation}

\subsection{The tilted field and the first particle lists}

Return to the two species and put $\mathbf1=(1,1)^t$.  The source in the
regulator is $im\int_TDa_T\,\mathbf1\cdot\chi$, since
$\widehat\lambda_0=D\widehat a$.  The curvature just computed satisfies
\[
\frac12\int_T(ma_T\mathbf1-i\chi)^tD(ma_T\mathbf1-i\chi)
+im\int_TDa_T\,\mathbf1\cdot\chi
=m^2\int_Ta_TDa_T-\frac12\int_T\chi^tD\chi.
\]
This is an identity for the common factor of each fixed current coefficient.
The odd real transforms of $a$ and $\lambda_0$ make them purely imaginary
on the real line; $D$ is real and positive.  Thus the remaining quadratic
factor is a real Gaussian tilt.  Including division by the original source
expectation gives
\begin{equation}\label{c7:source-removal}
\widetilde G=(I+\widehat GD)^{-1}\widehat G,\qquad
D_g=\prod_{n\ge1}\det(I+E(k)g(k))^{-1},\qquad
\kappa_T=\int_T a_TD(I+(\widehat G_d+\widehat G_c)D)a_T,
\end{equation}
with overall scalar $D_g e^{m^2\kappa_T}$ and $k=2\pi n/T$ in the product.
The determinant has power $-1$ because each positive frequency represents
the two real modes $k,-k$.  The products converge at fixed period: the zero
frequency is excluded and $gE$ decays exponentially at high frequency.
We henceforth let $\chi$ have this centered real tilted law.  This operation
is performed on the common factor before summing the new integrated series.

Combine \eqref{c7:oscillator-dresses} with the mean/winding field dresses:
\[
f_s(w)=f_s^{\rm osc}(w)-\alpha_pw-g_ww^2/2,\qquad
f_o(w)=f_o^{\rm osc}(w)-i(\alpha_p+g_ww)/2.
\]
Their multipliers give
\[
e^{f_s(v-i/2)-f_s(v+i/2)}=b(v)^{-1},\qquad
e^{f_o(v-i)+f_o(v+i)}=b(v).
\]
Define $H(v)=e^{f_s(v+i/2)-f_o(v-i)}$.  The nonzero-mode multiplier in its
exponent is $-E$ on $\mathsf A$, while the mean/winding difference is
$-\alpha_pv-g_wv^2/2+5g_w/8$.  Therefore
\begin{equation}\label{c7:H-dress}
\begin{gathered}
H(v)=B_T(v)^m
 \exp(2\pi m/T+\xi(v)-\alpha_pv-g_wv^2/2+5g_w/8),\\
B(v)=\frac{\cosh(\pi v/3)-\sqrt3/2}
           {\cosh(\pi v/3)+\sqrt3/2},\qquad
\widehat\xi(k)=iE(k)\widehat\chi(k)\quad(k\ne0).
\end{gathered}
\end{equation}
Here $B_T$ is the periodized product.  Indeed
\[
\widehat{\log B}(k)=-(2\pi/k)\sinh k/\cosh(3k/2)
=-E(k)\widehat a(k),\qquad \widehat{\log B}(0)=-2\pi.
\]
Removing this zero coefficient accounts for $2\pi m/T$ in
\eqref{c7:H-dress}.  The field $\xi$ is real on the real line: $E$ is real
and odd.  Its high-frequency covariance below shows that it is analytic on
$|\Im v|<3/2$.  The displayed exponential recurrences continue from their
defining strips.  For example, for $1/2<\Im v<3/2$, the upper value of $H$
can also be written as $e^{f_s(v-i/2)-f_o(v-i)}b(v)$; the first recurrence
provides that continuation.  The lower continuation is analogous.

We now specify the current pairs as particles.  A short or outer entry
$(a,\epsilon)$ means a field at $w=x+ia$ of sign $\epsilon$; an $H$ entry
$(a,r)$ means the full factor $H(x+ia)^r$.  All entries in one row belong to
one species.  The four rows $P,M,O_+,O_-$ come from roots $12,15,13,16$ in
$D$.  The rows $U,L$ come from the rightmost root and the cycled leftmost
root respectively.  The separate sign is their scalar coefficient sign.
\begin{equation}\label{j:particle-list}
\begin{array}{c|c|c|c|l}
 &\text{shorts}&\text{outers}&\text{sign}&H\text{-charges}\\\hline
P &(1/2,1)&(1,1)&+&(0,1)\\
M &(1/2,-1)&(1,1)&+&(0,1),(1,-1),(-1,-1)\\
O_+&\varnothing&(1,1),(0,1)&+&(0,1),(1,-1)\\
O_-&\varnothing&(1,1),(0,-1)&-&(0,1),(-1,-1)\\
U &(d,-1)&(1/2+d,1)&-&(3/2+d,-1)\\
L &(-d,-1)&(-1/2-d,-1)&+&(-3/2-d,-1)
\end{array}
\end{equation}
The internal order puts the height-$1$ outer first in each $D$ row, and the
short first in $U,L$.  To verify the table, a cycled $D$ pair has fields
$\Psi_3\Psi_{j-1}$ at $x+i$ and its coefficient at $x$.  The recurrences
give the four dresses
\[
H(x),\quad \frac{H(x)}{H(x+i)H(x-i)},\quad
\frac{H(x)}{H(x+i)},\quad -\frac{H(x)}{H(x-i)}.
\]
For $U$ at folded coordinate $v=x+i(1/2+d)$ the dress is $-H(v+i)^{-1}$;
for $L$ at its pre-cycling coordinate $v=x-i(1/2+d)$ it is $H(v-i)^{-1}$.
These are exactly the listed heights and signs.

At this stage fix $d=d_-<0$ with $|d_-|$ sufficiently small.  Starting from
the outgoing ordered strip in the preceding subsection, move the folded
$U$ coordinates to imaginary part $1/2+d_-$, the cycled $L$ coordinates to
$1/2-d_-$, and the cycled $D$ coordinates to $1$ by a limit from below.
Repeated pairs retain close separations until their analytic coincident
limits are taken.  These moves remain in the ordered strip.  The original
$b$ coefficients have no poles there by their $U_\delta$ products, and a
possible dual-pair singularity between $D$ pairs is cancelled by the zero
from their common first field.  In particular the extreme $H$ heights have
absolute value $3/2-|d_-|<3/2$.  The moves and the table thus hold for every
fixed coefficient before Gaussian integration.

Let $e_1,e_2$ be the species basis.  For a particle $j$ in species $\sigma$
define $q_j$ as its total short charge times $e_\sigma$.  Inspection of the
table shows that the total of its $H$-charges is the same vector.  The six
values are $(1,-1,0,0,-1,-1)e_\sigma$.  Consequently the trace constraints
\eqref{c7:charge-closure} imply, for every contributing term,
\begin{equation}\label{c7:gas-charge-closure}
\sum_jq_j=\boldsymbol\ell=(\ell_1,\ell_2)\in8\mathbb Z(1,-1).
\end{equation}
This constraint is retained in all sums below.

\subsection{Absolute integration on the first contours}

We first record the covariance used in the absolute estimate.  Write
$P_+$ for the common and $P_-$ for the difference species projection.  At nonzero
frequency the covariance of $\xi$ has symbol $2\pi\mathsf h(k)/k$, where
\begin{equation}\label{c7:xi-symbol}
\begin{gathered}
\mathsf h=gE^2(I+gE)^{-1},\qquad
\mathcal S=P_-+P_+/9,\qquad \mathcal L=I-\mathcal S=(8/9)P_+,\\
\mathsf h(k)/k=\mathcal S/k^2+c_0+O(k^2),\qquad
c_{0,-}=-17/12,\quad c_{0,+}=-25/108.
\end{gathered}
\end{equation}
For example, $E=1/k+7k/12+O(k^3)$,
$g_+=k/8-43k^3/96+O(k^5)$, and
$g_-=1/(2k)-17k/24+O(k^3)$ give these expansions.  With $t=e^{-k}$ the
exact channel symbols are
\[
\mathsf h_+=\frac{t^4(t+1)}{(1-t)(1+t^2)(1-t+t^2)(1+t+t^2)^2},\qquad
\mathsf h_-=\frac{t^3(4t^4-7t^3+10t^2-7t+4)}
 {(1-t)(1+t)(1+t^2)(1-t+t^2)^3}.
\]

Define $K^r$ on the real line by transform
$2\pi(\mathsf h/k-\mathcal S/k^2)$.  It is continuous, even, and
exponentially decaying.  More precisely, for $s>0$,
\begin{equation}\label{c7:residual-series}
K^r(s)=\sum_{j\ge1}P_j(s)e^{-\pi j s/6},
\end{equation}
where the matrix polynomials have uniformly bounded degree and coefficients
of at most polynomial growth.  The series also defines $K^r$ for $\Re s>0$.
To see this, the transform is removable at zero and its other poles lie on
the imaginary $i\pi/6$ grid with bounded order.  Shift Fourier inversion to
lines between successive grid points.  On those lines $\mathsf h$ is
$O(e^{-3|\Re k|})$ uniformly, while the subtracted rational term is
integrable; the residues give \eqref{c7:residual-series}.  The principal
parts of the rational functions of $e^{-k}$ are periodic along the grid,
which gives the stated coefficient bounds.

The form convenient at coincident horizontal coordinates is
\begin{equation}\label{c7:K0}
K_0(s)=K^r(0)+\int_{\mathbb R}(\cos(ks)-1)\mathsf h(k)\frac{dk}{k}.
\end{equation}
It equals $K^r(s)-\pi|s|\mathcal S$ for real $s$.  It is analytic on
$|\Im s|<3$.  Since the leading high-frequency term is $4P_-e^{-3k}$ for
$k>0$, its continuation to $|\Im s|<4$ has the sole singular part
\begin{equation}\label{c7:K0-singular}
-4P_-\log(1+s^2/9).
\end{equation}
In particular the exponentials between integer species charges are
meromorphic in that wider strip.  The zero-mean periodic covariance is
\begin{equation}\label{c7:periodic-K}
K_T(s)=\sum_{j\in\mathbb Z}K^r(s+jT)
 +\pi(s^2/T-|s|+T/6)\mathcal S-2\pi c_0/T,
\qquad -T<s<T.
\end{equation}
Fourier series proves this identity.  At complex shifts the term involving
$K^r(s)-\pi|s|\mathcal S$ is read as $K_0(s)$ and the other terms continue
analytically.

The following calculation concerns one finite list of $H$-charges.  Write
$r_i$ for its signed species vector, $w_i=x_i+ia_i$, and
\[
X=\sum_i r_iw_i,\qquad X^{(2)}=\sum_i r_iw_i^2.
\]
All squares in analytic formulas are bilinear.  Wick's formula for this
list is $\exp(\tfrac12\sum_{i,l}r_i^tK_T(w_l-w_i)r_l)$.  Under
\eqref{c7:gas-charge-closure}, $\mathcal S\boldsymbol\ell=\boldsymbol\ell$.
The quadratic and constant terms of \eqref{c7:periodic-K} therefore
contribute
\begin{equation}\label{c7:wick-quadratic}
\exp\!\left[
\frac\pi T(\boldsymbol\ell\cdot X^{(2)}-X^t\mathcal S X)
+\frac{\pi T}{12}|\boldsymbol\ell|^2
+\frac{17\pi}{12T}|\boldsymbol\ell|^2\right].
\end{equation}
The first term cancels the $-g_ww_i^2/2$ dresses.  The remaining winding
terms cancel exactly: \eqref{c7:modular-winding} cancels the term of order
$T$, and the coefficients of $\pi|\boldsymbol\ell|^2/T$ from Wick, the
$5g_w/8$ dresses, the modular scalar, and the normalized regulator weight
are respectively
\[
\frac{17}{12}+\frac{15}{12}-\frac8{12}-\frac{24}{12}=0.
\]
The common factor $e^{2\pi m/T}$ in $H$ also cancels, since
$\mathbf1\cdot\boldsymbol\ell=0$.

We apply this algebra first to an absolute moment on the $d_-$ contours.
Write $r_i=t_i e_{\sigma_i}$ with $t_i\in\mathbb R$.  Because each
$\xi_\sigma$ is real on the real line,
\[
\left|e^{\sum_i t_i\xi_{\sigma_i}(x_i+ia_i)}\right|
=e^{\sum_i t_i\Re\xi_{\sigma_i}(x_i+ia_i)}.
\]
Writing $r@w$ for a signed species charge $r$ at $w$, its expectation under
the tilted law is precisely the ordinary Wick exponential obtained by
replacing each signed charge by two signed halves,
\begin{equation}\label{c7:symmetrized-charges}
r_i@(x_i+ia_i)\quad\longmapsto\quad
\tfrac12r_i@(x_i+ia_i)+\tfrac12r_i@(x_i-ia_i).
\end{equation}
No sign of $t_i$ is changed.  This is a finite real Gaussian moment: the
largest height is $3/2-|d_-|$, so the high-frequency variance has a positive
exponential gap against $\mathsf h_-(k)=O(e^{-3k})$.  In channel $\varepsilon=\pm$,
the resulting height kernel is
$-\mathsf h_\varepsilon(k)\cosh(a_i k)\cosh(a_l k)$ in the negative logarithm of
the residual magnitude.  No continuation term is present.

The symmetrized list has the same total $\boldsymbol\ell$, and its two
moments are
\begin{equation}\label{c7:symmetrized-moments}
\widetilde X=\Re X=\sum_i r_ix_i,\qquad
\widetilde X^{(2)}=\Re X^{(2)}=\sum_i r_i(x_i^2-a_i^2).
\end{equation}
In particular the second moment retains $-a_i^2$.  Substituting these
moments in \eqref{c7:wick-quadratic} proves exactly the same winding
cancellations for the absolute moment.  Once those cancellations are made,
$\boldsymbol\ell$ remains fixed by \eqref{c7:gas-charge-closure}; it is not
an additional unrestricted summation index.

The heat sum is bounded before integration over constant modes.  Write
$\theta=\vartheta+i\tau$, $X=Y+iZ$, and
$\boldsymbol\alpha_{\mathbf p}=a_{\mathbf p}+i\tau\mathbf1$ with
$a_{\mathbf p}$ real.  For each heat summand,
\begin{equation}\label{c7:absolute-heat}
\left|e^{-T\boldsymbol\alpha_{\mathbf p}^2/(4\pi)
              -\boldsymbol\alpha_{\mathbf p}\cdot X}\right|
=e^{-Ta_{\mathbf p}^2/(4\pi)-a_{\mathbf p}\cdot Y}
 e^{T\tau^2|\mathbf1|^2/(4\pi)+\tau\mathbf1\cdot Z}.
\end{equation}
The second factor costs $e^{O_\theta(T+N)}$, where $N$ is particle count,
because the height lists are bounded.  Apply the triangle inequality to the
heat sum.  At fixed $k_0=p_1+p_2$, the opposite uniform mode and the remaining
integer sum unfold to $du/(2\pi)$ on $(u,-u)$.  The common Gaussian has
covariance $\pi P_+/(4T)$.  Its real completion, including
$e^{-\pi Y^t\mathcal S Y/T}$, gives
\begin{equation}\label{c7:real-heat-completion}
\frac{\sqrt{8/9}}{\sqrt{2T}}\sum_{k_0\in\mathbb Z}
 e^{-T\beta_\vartheta^t\mathcal L\beta_\vartheta/(4\pi)
       -\beta_\vartheta^t\mathcal L Y},
\qquad \beta_\vartheta=(\vartheta+\pi k_0)\mathbf1.
\end{equation}
Here $Y_+=\mathbf1\cdot Y/\sqrt2$ and
$Y_-=(1,-1)\cdot Y/\sqrt2$.  The opposite integral contributes
$(2T)^{-1/2}e^{\pi Y_-^2/T}$, and the common integral contributes
$\sqrt{8/9}e^{\pi Y_+^2/(9T)}$ as well as the displayed terms.  These two
exponentials cancel the corresponding parts of $-\pi Y^t\mathcal S Y/T$.
The two factors $\sqrt T$ from Poisson summation have been kept outside this
calculation.

For this real majorant, order the real particle coordinates and put
$Q(x)=\sum_{x_j<x}q_j$.  The real long-distance short and $H$ terms are
\[
\pi\sum_{u<j}(x_j-x_u)q_u^t\mathcal Lq_j.
\]
For each $k_0$ summand in \eqref{c7:real-heat-completion}, its exponent
together with these terms is the exponent of the real state square below,
with $p(x)=p_0-Q(x)$ and $p_0\in\mathbb Z^2$ having sum $k_0$:
\[
\exp\!\left[-\pi\int_T
 (p+(\vartheta/(2\pi))\mathbf1)^t\mathcal L
 (p+(\vartheta/(2\pi))\mathbf1)\,dx\right].
\]
Closure holds modulo $8\mathbb Z(1,-1)$ by
\eqref{c7:gas-charge-closure}.  For each such path, the base real product
consists of this state square, the short and outer local and residual
magnitudes, and the local and residual Wick factors from the complete
symmetrized $H$ lists.  It is the weight
$W^{\rm pre}_{d_-,\vartheta}(\mathcal Y,p)$ in
Theorem~\ref{j:absolute-stability}, with state angle
$\vartheta=\Re\theta$.  The explicit second factor of
\eqref{c7:absolute-heat} is kept outside this product and bounded separately
by $e^{C_\Omega(T+N)}$ on a compact set $\Omega$ of angles.  The height
phases have disappeared under the modulus, so this base product has no Berry
factor.  In particular no signed phase formula is applied to the symmetrized
list, whose $H$ height moment is zero.

At fixed $m,T,d_-$ the deterministic $B_T$ factors on these contours have a
finite bound per particle.  The pre-collision clause of
Theorem~\ref{j:absolute-stability}, after absorbing the heat factor
$e^{C_\Omega(T+N)}$ and the deterministic linear count cost,
therefore proves the required absolute summability on its domain of integer
periods for every finite $m$.  The allowed periods are independent of $m$;
only the absorbed count cost changes.  More explicitly, let $\nu$ index a
finite current monomial with specified species and roots, with the winding fixed
by \eqref{c7:gas-charge-closure}.  Let $I_{\nu,\mathbf p}$ be its full
conditional scalar coefficient after the common factor in
\eqref{c7:source-removal} is extracted.  It includes the exponential
factorial, modular field, pair, and image factors, table signs and $H$
dresses, the $\mathbf p$ heat factor, and the original winding probability.
Then
\begin{equation}\label{c7:pre-tonelli}
\sum_\nu\sum_{\mathbf p\in\mathbb Z^2}
\int_{\Gamma_-^\nu}\int d\mu_{\rm const}\,
\mathbb E_{\rm tilt}|I_{\nu,\mathbf p}|<\infty.
\end{equation}
Here $\Gamma_-^\nu$ is the product of the fixed $d_-$ contours, integrated
with product arclength, equivalently $\prod_jdx_j$ on the shifted horizontal
lines, and
$d\mu_{\rm const}$ is the original common-Gaussian and opposite-uniform
constant-mode law.  To see the final summation explicitly, if $n_j$ counts
particles in unit box $j$ and $p_{c,j}=(p_1+p_2)/\sqrt2$ is the state at
its left endpoint, stability bounds the magnitude by
$\exp(CT-c\sum_j(n_j^2+p_{c,j}^2))$, up to the absorbed linear costs.
For a finite label count $J$, every box contributes at most a sum of the
form
\[
\sum_{n\ge0}\frac{J^n}{n!}e^{-cn^2+A_{m,T,d_-}n}<\infty.
\]
The initial-state sum is Gaussian in $k_0=\sqrt2p_{c,0}$ and has eight
quotient classes for each $k_0$.  Dropping other configuration constraints
only enlarges this bound.  There is no further winding multiplicity.
This proves \eqref{c7:pre-tonelli}, locally uniformly for $\theta$ in compact
sets.  Its constants may depend on the fixed $d_-$; it asserts no uniform
estimate as $d_-\uparrow0$.

Tonelli and Fubini now pass the conditional current identities through the
Gaussian and constant-mode integrations.  This is the point at which the
regulator becomes the sum of the signed Wick-integrated terms on the first
contours.  The conditional trace calculation alone was not used for that
interchange.

\subsection{Wick integration and the collision residue}

For each signed term furnished by \eqref{c7:pre-tonelli}, use the original
integer $H$-charges in Wick's formula.  Equation~\eqref{c7:wick-quadratic}
and its winding cancellations now hold with complex $X,X^{(2)}$.
The signed constant-mode integration is the analytic version of
\eqref{c7:real-heat-completion}:
\begin{equation}\label{c7:signed-heat-completion}
\frac{\sqrt{8/9}}{\sqrt{2T}}\sum_{k_0\in\mathbb Z}
 \exp[-T\beta^t\mathcal L\beta/(4\pi)-\beta^t\mathcal L X],
\qquad \beta=(\theta+\pi k_0)\mathbf1.
\end{equation}
The two Poisson factors $\sqrt T$ remain outside.  For a fixed finite field
list the full heat sum converges normally on compact contour regions, with a
bound of the form $e^{-cTk_0^2+C_N|k_0|}$.  Retaining this sum gives the
horizontal periodicity inherited from the original current expression.
This periodic meromorphic expression is the one used in the contour moves;
an individual heat summand need not be periodic.

We make the local normalization explicit before calculating the residue.
For a basic type $v$, let $N_s(v),N_o(v)$ be its field counts, let $I_v$
count outer-before-short pairs inside its displayed internal order, and let
$\tau_v$ be its table sign.  The global outer-before-short sign from
\eqref{c7:klein-sign} can be assigned locally as
$(-1)^{I_v}i^{N_s(v)\bmod2}$.  Indeed a particle has the same parity of
short and outer counts.  If $N_{\rm odd}$ particles have odd short count, the
interparticle inversion exponent is $N_{\rm odd}(N_{\rm odd}-1)/2$.  The
trace constraint makes $N_{\rm odd}$ even in each species, and then
$(-1)^{N_{\rm odd}(N_{\rm odd}-1)/2}=i^{N_{\rm odd}}$.

Put $K_{\rm diag}(s)=(K_0(s))_{\sigma\sigma}$, the same for either species,
and $\kappa_j^{\rm loc}=K_{\rm diag}(ji)$ for $j=0,1,2$.  For each type $v$, denote its
ordered short, outer, and $H$ lists by $\mathscr S_v,\mathscr O_v,\mathscr H_v$.
The constant local factor of a basic particle is
\begin{equation}\label{c7:local-factor-rule}
\begin{split}
\gamma_v={}&\frac{a_*}{2\pi i}\tau_v(-1)^{I_v}i^{N_s(v)\bmod2}
 \left(\frac{2\pi}{a_*}\right)^{N_s(v)/2}
 \left(\frac{\pi}{4a_*}\right)^{N_o(v)/2}\\
&\times\prod_{i<l\text{ in }\mathscr S_v}
 [2\sinh(\pi i(a_l-a_i))]^{\epsilon_i\epsilon_l}
 \prod_{i<l\text{ in }\mathscr O_v}
 [\tanh(\pi i(a_l-a_i)/4)]^{\epsilon_i\epsilon_l}\\
&\times\exp\!\left(\tfrac12\sum_{i,l\text{ in }\mathscr H_v}
 r_ir_l K_{\rm diag}(i(a_l-a_i))\right).
\end{split}
\end{equation}
The factor uses $K_0$ within a particle; image factors remain outside it.
Since $a_* /(2\pi i)=-i/12$, the short and outer field norms are
$\sqrt{12}$ and $\sqrt{3/2}$.  Substitution gives
\begin{equation}\label{c7:basic-activities}
\begin{aligned}
\gamma_P=\gamma_U&=-\frac{\sqrt2}{4}e^{\kappa_0^{\rm loc}/2},
&\gamma_L&=\frac{\sqrt2}{4}e^{\kappa_0^{\rm loc}/2},\\
\gamma_M&=-\frac{\sqrt2}{4}e^{3\kappa_0^{\rm loc}/2-2\kappa_1^{\rm loc}+\kappa_2^{\rm loc}},
&\gamma_{O_+}=\gamma_{O_-}&=-\frac18e^{\kappa_0^{\rm loc}-\kappa_1^{\rm loc}}.
\end{aligned}
\end{equation}
For the $O$ rows the outer base at height difference $-1$ is $-i$; its
inverse for $O_-$ accounts for the equality of the two constants.  In
particular $P$ and $M$ have nonzero local factors.

Choose a fixed $d_+>0$ in the post-collision stability range, with
$|d_-|<d_+$; both are small enough for the strips used below.  In each finite
term move the $U$ short-center contours upward first, from height $d_-$ to
$d_+$.  They cross the $L$ short-center contours, initially at height
$-d_-$.  Then move the remaining $L$ contours downward to $-d_+$ and move
any fused centers to the real line.  A lower horizontal contour minus an
upper one is the positive boundary of the intervening cylinder strip.
Thus a crossed pole contributes $+2\pi i$ times its residue in the $U$
center.

For a same-species $U,L$ pair write $\delta=z_U-z_L$, where $z_U,z_L$ are
their short centers.  Their ordered fields are
\[
(z_U,-)_{s},\quad(z_U+i/2,+)_{o},\quad
(z_L,-)_{s},\quad(z_L-i/2,-)_{o}.
\]
Their $H$-charges are $-1$ at $z_U+3i/2$ and $z_L-3i/2$.
The joint table sign is negative and the ordered union has one
outer-before-short inversion, so those two signs cancel.  Its two shorts
have even parity.  The current and field-norm factors together are $-1/8$.
Therefore the full local germ is
\begin{equation}\label{c7:fusion-germ}
-\frac18\,2\sinh(-\pi\delta)\,
\coth\!\left(-\frac{\pi(\delta+i)}4\right)
\exp[K_{\rm diag}(0)+K_{\rm diag}(\delta+3i)].
\end{equation}
Define $K_{\rm diag}^{\rm reg}(s)=K_{\rm diag}(s)+2\log(1+s^2/9)$ by continuation from the inner
strip.  By \eqref{c7:K0-singular} it is regular at $3i$, and
\[
e^{K_{\rm diag}(\delta+3i)}
=\frac{81e^{K_{\rm diag}^{\rm reg}(\delta+3i)}}{\delta^2(\delta+6i)^2}
\sim-\frac9{4\delta^2}e^{K_{\rm diag}^{\rm reg}(3i)}.
\]
Using $2\sinh(-\pi\delta)\sim-2\pi\delta$ and
$\coth(-i\pi/4)=i$, the residue of \eqref{c7:fusion-germ} in $z_U$ is
\[
-\frac{9\pi i}{16}\exp[K_{\rm diag}(0)+K_{\rm diag}^{\rm reg}(3i)].
\]
The fused local activity is hence
\begin{equation}\label{c7:fused-activity}
\gamma_{F_{\rm pair}}=\frac{9\pi^2}{8}\exp[K_{\rm diag}(0)+K_{\rm diag}^{\rm reg}(3i)]>0.
\end{equation}
The singular power is the integer $-2$, so it has no branch sign; also
$d\delta/dz_U=1$.  The regular value is real by continuation along $s=iy$
with $y<3$, proving the last inequality.

Call this fused particle $F_{\rm pair}$.  Its internal order is $U$ then
$L$, and its complete lists, at its real center $x$, are
\begin{equation}\label{c7:fused-list}
\begin{array}{c|c|c}
\text{shorts}&\text{outers}&H\text{-charges}\\\hline
(0,-1),(0,-1)&(1/2,1),(-1/2,-1)&(3/2,-1),(-3/2,-1).
\end{array}
\end{equation}
These are the lists $\mathscr S_{F_{\rm pair}},\mathscr O_{F_{\rm pair}},
\mathscr H_{F_{\rm pair}}$.  The local activity has already taken the residue
of the entire factor \eqref{c7:fusion-germ}; no internal pair product is
evaluated at its coincident short centers.
External interactions and nonzero images simply evaluate at equality and
combine into these union lists.  The constant \eqref{c7:fused-activity} is
independent of $T,m$ and the state.  In particular the local germ uses
$K_0$, before any long-distance separation; the corresponding constant
from $K^r$ alone would have an extra phase.

Here are all possible divisor crossings in these moves.  From
\eqref{c7:K0-singular}, two extreme $H$-charges of sign $-1$ in the same
species give a double pole at height difference $3$; in different species
they give a double zero.  The two negative shorts in a $U,L$ pair give one
zero, leaving exactly the simple same-species pole used above.  Against a
fused particle, $U$ or $L$ has two short zeros and at most one double
extreme-height pole.  Two fused particles have four short zeros and at most
two such double poles.  Thus these later encounters are regular and cannot
create a second fusion of the same particle.  A same-species $U,U$ or $L,L$
coincidence has like-charge short and outer zeros and equal $H$ heights; a
cross-species coincidence has only the regular equal-height $H$ factor.
Neither has a pole.

Among $D$ rows, a $P,M$ coincidence has one inverse short factor, cancelled
by the zero of their common height-$1$ outer field.  An $O_+,O_-$
coincidence has one inverse height-$0$ outer factor, cancelled by the same
common outer zero.  The other $D$ coincidences have no uncancelled inverse
factor.  Their $H$ height differences are at most $2$.  Mixed $D$ and
$U,L,F_{\rm pair}$ $H$ differences are at most $5/2+d_+<3$; the other
outer differences remain below $2$, and no other short divisor is crossed.
Finally
\begin{equation}\label{c7:B-recurrences}
B(v+i)B(v-i)=B(v),\qquad B(v+3i)=B(v)^{-1}
\end{equation}
gives the simplified $D$ factors $B(x),1,B(x-i),B(x+i)$ and makes the $U,L$
factors cancel at fusion.  The zeros of $B$ are at $i(\pm1/2+6n)$ and its
poles are at $i(\pm5/2+6n)$, $n\in\mathbb Z$.  The moving inverse $U,L$
arguments have heights in $(1,2)$ and $(-2,-1)$, respectively, while the
$D$ arguments have heights $0,\pm1$.  Thus these factors meet no divisor
during the collision moves.  The same statements hold for their periodized
products.  They account for every pair of rows in
\eqref{j:particle-list} and \eqref{c7:fused-list}.

At fixed particle number the preceding statements are identities of
periodic meromorphic contour integrals.  Each $U$ may fuse with at most one
$L$, and fusions occur only within a species.  For $u,l$ original particles
and $r$ disjoint pairs the multiplicity is
\begin{equation}\label{c7:fusion-factorial}
\frac{\binom ur\binom lr r!}{u!\,l!}
=\frac1{(u-r)!(l-r)!r!}.
\end{equation}
There is no additional $1/2$, because $U,L$ are distinct and the residue is
always in the $U$ center.  Reordering whole bilinears has even parity, and
each matching gives the same ordered union and sign.  Thus the formula also
counts the absolute values of the individual residue contributions.

\subsection{The post-collision data and endpoint summability}

From now on write $d=d_+>0$.  For a type $v$ define its scalar charge and
height moments by
\[
q_v=\sum_{(a,\epsilon)\in\mathscr S_v}\epsilon
    =\sum_{(a,r)\in\mathscr H_v}r,\qquad
s_v=\sum_{(a,\epsilon)\in\mathscr S_v}a\epsilon,\qquad
b_v=\sum_{(a,r)\in\mathscr H_v}ar.
\]
For a particle of this type in species $\sigma$, the vectors used below are
$q_j=q_ve_\sigma$, $s_j=s_ve_\sigma$, and $b_j=b_ve_\sigma$.
Let $\omega_v$ be the product of its $B$ factors with the signed $H$ powers.
The recurrences \eqref{c7:B-recurrences} give the complete table
\begin{equation}\label{c7:particle-data}
\begin{array}{c|r|r|r|l}
v&q_v&s_v&b_v&\omega_v(x)\\\hline
P&1&1/2&0&B(x)\\
M&-1&-1/2&0&1\\
O_+&0&0&-1&B(x-i)\\
O_-&0&0&1&B(x+i)\\
U&-1&-d&-3/2-d&B(x+i(3/2+d))^{-1}\\
L&-1&d&3/2+d&B(x-i(3/2+d))^{-1}\\
F_{\rm pair}&-2&0&0&1.
\end{array}
\end{equation}
The periodized product is denoted $\omega_{v,T}$, with value $1$ in the two
constant rows.  The original deterministic factor is exactly
$\omega_{v,T}(x)^m$.  These simplified expressions also specify the value
when separate factors in a ratio have cancelling zeros and poles.  The
physical fused particle has charge $-2e_\sigma$; the normalized half-fused
vector used in the matrix proof of stability is only a coordinate in that
proof.

For positive real separation $x$ and species-labelled types $u,v$, define
the forward residual factor using their complete lists by
\begin{equation}\label{c7:forward-residual}
\begin{split}
\mathscr R_{uv}(x)={}&
\prod_{\substack{(a,\epsilon)\in\mathscr S_u\\
                 (b,\epsilon')\in\mathscr S_v\\
                 \sigma_u=\sigma_v}}
 [1-e^{-2\pi(x+i(b-a))}]^{\epsilon\epsilon'}\\
&\times
\prod_{\substack{(a,\epsilon)\in\mathscr O_u\\
                 (b,\epsilon')\in\mathscr O_v\\
                 \sigma_u=\sigma_v}}
 [\tanh(\pi(x+i(b-a))/4)]^{\epsilon\epsilon'}\\
&\times\exp\!\left[
\sum_{\substack{(a,r)\in\mathscr H_u\\(b,r')\in\mathscr H_v}}
 rr'e_{\sigma_u}^tK^r(x+i(b-a))e_{\sigma_v}\right].
\end{split}
\end{equation}
An empty product is one.  The series \eqref{c7:residual-series} and the
logarithmic short and outer expansions make its log exponentially decaying
at large $x$.  On a periodic configuration, its forward product takes each
particle in one period as source and every strictly later particle in the
periodized sequence as target.  In particular it includes all nonzero
self-images.  Only the true diagonal at separation zero is omitted.

We next justify summing the finite contour identities.  For each summand of
\eqref{c7:signed-heat-completion}, write $\beta=\beta_\vartheta+i\tau\mathbf1$
and $X=Y+iZ$.  Its modulus is the real heat summand with $\vartheta$ and $Y$,
times
$\exp[T\tau^2\mathbf1^t\mathcal L\mathbf1/(4\pi)+\tau\mathbf1^t\mathcal L Z]$.
On a compact $\theta$ set the latter costs $e^{O_\theta(T+N)}$.
For distinct sorted particles let $A_{uv}(x)$ be the direct pair factor
obtained from $\mathscr R_{uv}(x)$ by replacing each short base by
$2\sinh(\pi(x+i(b-a)))$ and each $K^r$ by $K_0$.  The relation
$K_0(s)=K^r(s)-\pi s\mathcal S$ on $\Re s>0$ gives directly
\[
|A_{uv}(x)|=e^{\pi xq_u^t\mathcal Lq_v}|\mathscr R_{uv}(x)|.
\]
All height terms in the ratio are phases, also after continuation across
height $3$ for the integer $H$-charges.  The images are already residual
factors.  The real long-distance part therefore completes into the common
state square used above, and the remaining pair magnitudes are exactly
$|\mathscr R_{uv}|$ from the actual integer height lists.  The local magnitudes are those of
\eqref{c7:basic-activities} and \eqref{c7:fused-activity}.  This is exactly
the post-collision real pair/state/local weight estimated by
Theorem~\ref{j:absolute-stability}.  The bounded imaginary-angle costs are
linear in the count and are absorbed there.

The deterministic factors require no added count cost on these contours.
Writing $z=\cosh(\pi(x+ia)/3)$ and $r=\sqrt3/2$ shows
$|z-r|<|z+r|$ precisely when $\cos(\pi a/3)>0$.  Hence $|B(x+ia)|<1$ for
$|a|<3/2$ and $|B(x+ia)|>1$ just beyond those lines.  Every nonconstant row
of \eqref{c7:particle-data} has $|\omega_v(x)|<1$, and this remains true
image by image.  The post-collision stability bound and the same count and
state summation used in \eqref{c7:pre-tonelli} therefore prove absolute
summability at this fixed $d_+$ for all supported periods.

Each finite term has already been deformed with its residues.  Equation
\eqref{c7:fusion-factorial} shows that the sum of their absolute residue
contributions is bounded by the ordinary post-collision factorial series,
so the two endpoint bounds permit summation of all those identities.  No
estimate uniform between $d_-$ and $d_+$ is required.  In particular the
raw Gaussian field was evaluated only on the first contours; the contour
crossing used its Wick-integrated meromorphic functions.

\subsection{The exact horizontal gas}

It remains to express each signed post-collision term in variables local to
the horizontal direction.  Order its distinct real centers in one period.
For a positive separation $s=x+i(b-a)$, the full short factor splits as
$2\sinh\pi s=e^{\pi s}(1-e^{-2\pi s})$, and
$K_0(s)=K^r(s)-\pi s\mathcal S$ on $\Re s>0$.  Keep the full $K_0$ within a
particle.  For distinct particles $u<j$, summing the separated linear terms
from their lists gives
\[
\pi\sum_{u<j}(x_j-x_u)q_u^t\mathcal Lq_j
+i\pi\sum_j(s_j-\mathcal S b_j)\cdot
     (Q(x_j-)+Q(x_j+)-\boldsymbol\ell),
\qquad Q(x)=\sum_{x_j<x}q_j.
\]
The residual factors, including all images, are precisely
\eqref{c7:forward-residual}.

Choose $p_0\in\mathbb Z^2$ with $p_{0,1}+p_{0,2}=k_0$ and put
$p(x)=p_0-Q(x)$.  The real-coordinate linear term above and the heat and
mean terms of \eqref{c7:signed-heat-completion} complete to
\[
-\pi\int_T(p+(\theta/(2\pi))\mathbf1)^t\mathcal L
                 (p+(\theta/(2\pi))\mathbf1)\,dx.
\]
This follows by expanding the square, using
$\int_TQ\,dx=T\boldsymbol\ell-\sum_jq_jx_j$ and
$\mathcal L\boldsymbol\ell=0$.  The imaginary terms become
\[
-i\pi\sum_j(s_j-\mathcal S b_j)\cdot(p(x_j-)+p(x_j+))
-i\theta\mathbf1^t\mathcal L\sum_j b_j.
\]
To check the discarded phase exactly, \eqref{c7:particle-data} gives
\[
s_v-b_v=(1/2,-1/2,1,-1,3/2,-3/2,0)
\quad\text{for }P,M,O_+,O_-,U,L,F_{\rm pair}.
\]
It is congruent to $q_v/2$ modulo integers, so its total is integral in
each species by \eqref{c7:gas-charge-closure}.  The phase left by the
substitution is
\[
2\pi p_0\cdot\sum_j(s_j-b_j)
-\pi\boldsymbol\ell\cdot\sum_j(s_j-b_j),
\]
a multiple of $2\pi$.  Shifting $p_0$ by $8(1,-1)$ leaves each local weight
unchanged.  For a fixed sum $k_0$, the eight classes modulo this shift have
the same full-tour weight: shifting between them by $(1,-1)$ changes the
product phase by $2\pi$ times the same integral total.  This explains the
factor $1/8$ in the normalization below.

This computation also checks the fused local phase.  If $U$ precedes $L$
and $s=z_L-z_U=x-2id$, their full short and $H$ factors divided by their
residual factors give
\[
e^{\pi s-\pi\mathcal S_{\sigma\sigma}(s-3i)}
=e^{\pi\mathcal L_{\sigma\sigma}x}
 e^{i\pi(3\mathcal S_{\sigma\sigma}-2d\mathcal L_{\sigma\sigma})}.
\]
The first factor is their pair term in the state square.  The second is the
product of their local phases for successive states
$p,p+e_\sigma,p+2e_\sigma$; their $\theta$ phases cancel since $b_U+b_L=0$.
At collision this phase is already part of the full $K_0$ germ
\eqref{c7:fusion-germ}.  The fused row has $q=-2$, $s=b=0$ and therefore
requires no further local phase.

We can now define the post-collision gas.  Let
\[
\mathcal V=\{1,2\}\times\{P,M,O_+,O_-,U,L,F_{\rm pair}\}.
\]
A configuration $\mathcal Y=((x_j,v_j))_{j=1}^N$ consists of a finite number
of particles with $0\le x_j<T$ and $v_j\in\mathcal V$; coincident real
positions form a null set and are discarded.  For
$[p_0]\in\mathbb Z^2/8\mathbb Z(1,-1)$ let $p$ have the jumps $-q_j$ in
increasing order, using a consistent lift during the period, and impose
$p(T)=p(0)$ in the quotient.  Define
\begin{equation}\label{c7:state-weight}
\begin{split}
\mathcal Q_\theta(\mathcal Y,p)={}&
\exp\!\left[-\pi\int_T
 (p+(\theta/(2\pi))\mathbf1)^t\mathcal L
 (p+(\theta/(2\pi))\mathbf1)\,dx\right]\\
&\times\prod_j\exp\!\left[
-i\pi(s_j-\mathcal S b_j)\cdot(p(x_j-)+p(x_j+))
-i\theta\mathbf1^t\mathcal L b_j\right].
\end{split}
\end{equation}
The weight with the deterministic site factors omitted is
\begin{equation}\label{c7:horizontal-weight}
W_{d,\theta}(\mathcal Y,p)=\mathcal Q_\theta(\mathcal Y,p)
 \prod_j\gamma_{v_j}
 \prod_{\substack{1\le j,k\le N,\ n\in\mathbb Z\\x_k+nT>x_j}}
 \mathscr R_{v_jv_k}(x_k+nT-x_j).
\end{equation}
Thus $j=k,n=0$ is excluded and every
$j=k,n\ge1$ is included.  The constants $\gamma_v$ use
\eqref{c7:basic-activities} or \eqref{c7:fused-activity}, independently of
the species.  For the empty configuration the particle and residual products
are one, while the state heat factor remains.
For distinct positions the factors in \eqref{c7:horizontal-weight} are
continuous and nonzero: the local activities are nonzero, the short and
outer residual bases have no zero or pole at positive real separation, and
the remaining factors are exponentials.

For counts $(n_v)_{v\in\mathcal V}$ with $N=\sum_vn_v$, use any fixed
labeling of those particles in the integral and set
\begin{equation}\label{c7:gas-definition}
Z_{m,T}=\sum_{(n_v)\in\mathbb Z_{\ge0}^{\mathcal V}}
 \frac1{\prod_v n_v!}
 \sum_{[p_0]\in\mathbb Z^2/8\mathbb Z(1,-1)}
 \int_{[0,T)^N}\!\mathbf1_{\{p(T)=p(0)\}}
 W_{d,\theta}(\mathcal Y,p)
 \prod_{j=1}^N\omega_{v_j,T}(x_j)^m\,\prod_{j=1}^Ndx_j.
\end{equation}
The indicator uses quotient closure.  All definitions are independent of
the chosen lift by the preceding phase check.  The absolute endpoint bound
proves convergence of this series for the periods in the next proposition,
locally uniformly for $\theta$ in compact sets.  The local factors use
$K_0$ and the forward product uses $K^r$, so no complex convention for
$|s|$ is needed inside a particle.

\begin{proposition}[Exact gas representation]\label{c7:gas-identity}
Fix $d_-<0<d$ in the respective nonzero displacement ranges of
Theorem~\ref{j:absolute-stability}, with $0<|d_-|<d<1/2$.
Fix also a compact set $\Omega\subset\mathbb C$ of values of $\theta$.
There is an integer $T_0$, independent of $m$, such that for every integer
$T\ge T_0$ and every finite integer $m\ge0$, the gas
\eqref{c7:gas-definition} converges absolutely, uniformly for
$\theta\in\Omega$ at that $m,T$, and satisfies
\begin{equation}\label{c7:regulator-gas}
\mathfrak C_{m,T}=c(T)e^{m^2\kappa_T}Z_{m,T},\qquad
c(T)=\frac{\sqrt{8/9}}{8\sqrt{2T}\,D_\ell}
 \left(\frac{\sqrt T D_W}{Z_\Gamma}\right)^2D_g,
\qquad D_\ell=\sum_{j\in8\mathbb Z}e^{-4\pi j^2/T}.
\end{equation}
In particular it holds along $T=bn$ for any fixed integer block length
$b\ge T_0$ and every positive integer $n$.  This common period domain allows
$T$ to remain fixed while $m$ varies through finite integers.  The absolute
bounds may depend on $m,T$ and the fixed displacements; no bound uniform in
$m$ or limiting assertion as $d_-\uparrow0$ or $d\downarrow0$ is made.
\end{proposition}

\begin{proof}
The conditional trace and oscillator calculation gives each fixed current
coefficient, with two empty oscillator factors $D_W/Z_\Gamma$ and
two Poisson factors $\sqrt T$.  Source removal gives
$D_g e^{m^2\kappa_T}$.  Equation~\eqref{c7:pre-tonelli} justifies integrating
and summing these coefficients on the first contours.  Their signed Wick
and constant-mode integrations give \eqref{c7:signed-heat-completion}, with
factor $\sqrt{8/9}/\sqrt{2T}$, and cancel all winding-dependent exponentials,
leaving $D_\ell^{-1}$.

The finite periodic contour identities give precisely the basic and fused
terms with constants \eqref{c7:basic-activities}--\eqref{c7:fused-activity}.
Their endpoint absolute bounds and \eqref{c7:fusion-factorial} justify
summing them.  The signed algebra leading to
\eqref{c7:state-weight}--\eqref{c7:gas-definition} then identifies that sum
with $Z_{m,T}$.  The eight equal full-tour representatives supply the factor
$1/8$.  Multiplying the displayed constants proves
\eqref{c7:regulator-gas} on the supported periods.
\end{proof}

The table now exposes the properties used by the following sections.  There
are finitely many types with bounded integer jumps; the nonzero $P,M$
activities permit both signs of each species jump.  The two types whose
site multipliers are exactly one, $M,F_{\rm pair}$, have strictly negative
common charge.  Every other site multiplier has modulus below one on the
post-collision real line.  The forward residual log has the exponentially
decaying finite-degree expansion inherited from
\eqref{c7:residual-series}.  Theorem~\ref{j:absolute-stability} controls the
absolute closed-sequence weights, and Section~\ref{c7:compact} uses these
properties to build the horizontal operator.

\subsection{The scalar rates}

We finish by calculating the two scalar limits used in that operator
argument.  In terms of $t=e^{-k}$, the eigenvalues of $I+Eg$ are
\[
\frac{(1-t^6)^2}{(1-t^4)(1-t^8)},\qquad
\frac{(1+t^3)^4}{(1-t^4)(1-t^8)}.
\]
Riemann sums have only logarithmic endpoint singularities.  Using
\[
\int_0^\infty\log(1-e^{-ak})\,dk=-\frac{\pi^2}{6a},\qquad
\int_0^\infty\log(1+e^{-ak})\,dk=\frac{\pi^2}{12a},
\]
gives the rates $-13\pi/144$ for $D_g$ and $11\pi/144$ for each
$D_W/Z_\Gamma$.  The Gaussian sum has $D_\ell=\Theta(\sqrt T)$ by comparison
with its integral, so it and the powers of $T$ are subexponential.  Consequently
\begin{equation}\label{c7:c-rate}
\lim_{T\to\infty}\frac{\log c(T)}{T}=\frac\pi{16}.
\end{equation}

The symbols in \eqref{c7:source-removal}, with
$\widehat a(-k)=-\widehat a(k)$, similarly give
\begin{equation}\label{c7:kappa-rate}
\lim_{T\to\infty}\kappa_T
=-2\int_0^\infty
 \frac{t(1-t)(1-t^3)^2}{(1+t^2)(1+t^3)(1+t^4)}\frac{dk}{k}
=2\log\frac{2(2+\sqrt2)}9.
\end{equation}
Here is the evaluation of the integral.  Multiplying its rational function
of $t$ by $1-t^{24}$ gives coefficients $d_j=d_{24-j}$ for $1\le j\le23$,
whose first twelve values are
\[
(1,-1,-1,-2,3,3,1,-4,-3,-1,3,2).
\]
Their total after reflection is zero.  Let $R(t)$ denote the rational
function in the integrand, so $(1-t^{24})R(t)=\sum_jd_jt^j$.  The grouped
partial sum through period $N$ is
$R(e^{-k})(1-e^{-24(N+1)k})$.  After division by $k$ it is dominated by
$R(e^{-k})/k$, which is integrable because $R(e^{-k})=O(k^3)$ at zero and
$O(e^{-k})$ at infinity.  Dominated convergence therefore justifies the
period expansion.  Frullani's log difference on each block gives its
integral as $-\sum_jd_j\log(24n+j)$.  Pair $j$ with $24-j$,
using exponent $d_{12}/2$ at $j=12$.  The normalized product
\[
\frac{\prod_{n=0}^N(n+a)(n+1-a)}{N!(N+1)!}
\longrightarrow a\prod_{n\ge1}(1-a^2/n^2)=\frac{\sin\pi a}{\pi}
\]
then gives the exponentiated sum of the paired logarithms as
$\prod_{j=1}^{11}\sin(\pi j/24)^{d_j}$.  The sine product itself follows,
for example, by taking the limit in finite odd sine multiplication
\[
\frac{\sin x}{N'\sin(x/N')}=
\prod_{j=1}^{(N'-1)/2}
 \left(1-\frac{\sin^2(x/N')}{\sin^2(\pi j/N')}\right),
\qquad N'\text{ odd}.
\]
Writing $s_j=\sin(\pi j/24)$ and pairing complementary small angles reduces
the required product to
\[
\frac{(s_{11}s_5)^2}{s_9^2s_8^4}=
\frac{2(2+\sqrt2)}9,
\]
which proves \eqref{c7:kappa-rate}.  Proposition~\ref{c7:gas-identity}
combined with the fixed-site regulator limit of Section~\ref{c7:contour}
therefore supplies the period limit needed later, always at fixed $m$.
